% ============================================================
%  commands.tex — Custom macros for Econ PhD lecture notes
%  Last updated: April 2026
% ============================================================

% ------------------------------------------------------------
%  1. Packages (loaded here to keep main .tex clean)
% ------------------------------------------------------------
\usepackage{amsmath}

% ------------------------------------------------------------
%  2. Environment shortcuts
% ------------------------------------------------------------
\newcommand{\ba}{\begin{aligned}}
\newcommand{\ea}{\end{aligned}}
\newcommand{\bc}{\begin{cases}}
\newcommand{\ec}{\end{cases}}

% ------------------------------------------------------------
%  3. Blackboard bold (number sets)
% ------------------------------------------------------------
\newcommand{\RR}{\mathbb{R}}
\newcommand{\QQ}{\mathbb{Q}}
\newcommand{\ZZ}{\mathbb{Z}}
\newcommand{\CC}{\mathbb{C}}
\newcommand{\NN}{\mathbb{N}}
\newcommand{\FF}{\mathbb{F}}

% ------------------------------------------------------------
%  4. Calligraphic letters (collections, spaces, etc.)
% ------------------------------------------------------------
\renewcommand{\P}{\mathcal{P}}
\renewcommand{\S}{\mathcal{S}}   % overwrites §
\newcommand{\T}{\mathcal{T}}
\newcommand{\A}{\mathcal{A}}
\newcommand{\B}{\mathcal{B}}
\newcommand{\C}{\mathcal{C}}
\renewcommand{\L}{\mathcal{L}}

% ------------------------------------------------------------
%  5. Bold vectors  (single-letter shortcuts)
%     For ad-hoc vectors, use \vb{symbol} instead.
% ------------------------------------------------------------
\newcommand{\vb}[1]{\mathbf{#1}}          % general bold vector
\newcommand{\x}{\mathbf{x}}
\newcommand{\y}{\mathbf{y}}
\newcommand{\z}{\mathbf{z}}
\newcommand{\w}{\mathbf{w}}
\newcommand{\e}{\mathbf{e}}
\renewcommand{\b}{\mathbf{b}}             % overwrites bar-under accent
\newcommand{\p}{\mathbf{p}}
\newcommand{\q}{\mathbf{q}}
\newcommand{\0}{\mathbf{0}}
\newcommand{\1}{\mathbf{1}}
\newcommand{\I}{\mathbf{I}}
\newcommand{\X}{\mathbf{X}}

% ------------------------------------------------------------
%  6. Greek letter shortcuts
% ------------------------------------------------------------
\newcommand{\ve}{\varepsilon}
\newcommand{\s}{\sigma}
\renewcommand{\l}{\lambda}                % overwrites Polish ł

% ------------------------------------------------------------
%  7. Delimiters (auto-sizing)
% ------------------------------------------------------------
\newcommand{\br}[1]{\left[#1\right]}                 % brackets
\newcommand{\pr}[1]{\left(#1\right)}                  % parentheses
\renewcommand{\brace}[1]{\left\{#1\right\}}           % braces
\newcommand{\abs}[1]{\left|#1\right|}                  % absolute value
\newcommand{\norm}[2][1]{\|{#2}\|_{#1}}               % norm
\newcommand{\floor}[1]{\left\lfloor #1 \right\rfloor} % floor
\newcommand{\ceil}[1]{\left\lceil #1 \right\rceil}    % ceiling

% ------------------------------------------------------------
%  8. Operators (proper math-operator spacing)
% ------------------------------------------------------------
\DeclareMathOperator*{\argmax}{arg\,max}
\DeclareMathOperator*{\argmin}{arg\,min}
\DeclareMathOperator{\plim}{plim}
\DeclareMathOperator{\supp}{supp}
\DeclareMathOperator{\mnull}{null}
\DeclareMathOperator{\mspan}{span}

% ------------------------------------------------------------
%  9. Probability, expectation, variance
%     Usage:  \Pr{A \given B},  \E{X \given Y},  \Var{X}
% ------------------------------------------------------------
\newcommand{\given}{\;\middle|\;}                     % conditional bar
\renewcommand{\Pr}[1]{\operatorname{Pr}\!\left(#1\right)}
\newcommand{\E}[2][]{\mathbb{E}_{#1}\!\pr{#2}}
\newcommand{\Var}[2][]{\operatorname{Var}_{#1}\!\br{#2}}
\newcommand{\Cov}[2]{\operatorname{Cov}\!\br{#1,#2}}

% ------------------------------------------------------------
%  10. Convergence and asymptotics
% ------------------------------------------------------------
\newcommand{\dto}{\overset{d}{\to}}                   % convergence in distribution
\newcommand{\pto}{\overset{p}{\to}}                   % convergence in probability
\newcommand{\asim}{\overset{\text{a}}{\sim}}           % asymptotically distributed

% ------------------------------------------------------------
%  11. Calculus and analysis
% ------------------------------------------------------------
\renewcommand{\d}{\text{d}}                           % upright differential d
\newcommand{\dd}[2]{\frac{\d #1}{\d #2}}              % total derivative
\newcommand{\pd}[2]{\frac{\partial #1}{\partial #2}}  % partial derivative
\newcommand{\inv}[1]{{#1}^{-1}}                       % inverse

% ------------------------------------------------------------
%  12. Relations and logic
% ------------------------------------------------------------
\newcommand{\se}{\subseteq}
\newcommand{\bs}{\setminus}                           % set difference (use \bs not \c)
\newcommand{\st}{\text{ s.t. }}
\newcommand{\mif}{\text{ if }}

% ------------------------------------------------------------
%  13. Summation / union / intersection shortcuts
% ------------------------------------------------------------
\newcommand{\add}[1][i]{\sum_{#1 = 1}^n}
\newcommand{\ply}[1][i]{\prod_{#1 = 1}^n}
\newcommand{\sinf}[1][k]{\sum_{{#1} = 1}^{\infty}}
\newcommand{\uinf}[1][k]{\bigcup_{{#1} = 1}^{\infty}}
\newcommand{\iinf}[1][k]{\bigcap_{{#1} = 1}^{\infty}}

% ------------------------------------------------------------
%  14. Game theory / mechanism design
% ------------------------------------------------------------
\newcommand{\pref}{\succsim}                          % weak preference
\newcommand{\meet}{\wedge}                            % lattice meet
\newcommand{\join}{\vee}                              % lattice join

% ------------------------------------------------------------
%  15. Complex analysis (keep if needed)
% ------------------------------------------------------------
\renewcommand{\Re}{\operatorname{Re}}
\renewcommand{\Im}{\operatorname{Im}}

% ------------------------------------------------------------
%  16. Draft utilities
% ------------------------------------------------------------
\newif\ifdraft
\drafttrue   % set \draftfalse to hide all TODOs

\newcommand{\todo}[1]{%
  \ifdraft{\textcolor{red}{\textbf{[TODO: #1]}}}{}\fi
}

% ------------------------------------------------------------
%  17. Highlighting (for lecture notes)
% ------------------------------------------------------------
\newcommand{\match}[1]{\textcolor{ForestGreen}{$\boldsymbol{#1}$}}
